\section{Audit of the inherited weak-arithmetic proof}
\label{sec:inherited-audit}

The inherited theorem used above comes from the latest incoming
\path{pa194_ax11h_research_2026-11-14.zip}, identified by its exact
hash in Section~\ref{sec:new-relative-ipra}.  Its written iPRA
proof is stronger than the older placed packet with the same basename.
The archive expressly states that it supplies no new end-to-end Lean
or Rocq build.  Thus the incomplete proof-assistant status of the
repository and the existence of a later written mathematical argument
are compatible statements.

The audit inspected the physical, source-command, command-preservation,
premise-exchange, Foundation-adapter, term-leaf, forward-calibration and
final-equivalence arguments.  No concrete mathematical gap was found in
the inspected chain.  Three matters deserve explicit attention.

First, the physical proof is a symbolic induction proof with a finite
binding to each literal table.  Its configurations retain all evolving
register, gap and spare-list data.  Its exact primitive recursive clocks
and positive segment durations give both directions of halting
transport.  Complete blank initialization, including blank exteriors,
is verified separately.  A finite successful prefix alone would not
give this result.

Second, command-compiler correctness for arbitrary base stacks is proved
by bounded Boolean induction using explicit input envelopes.  The
forward PA compiler uses finite-demand query trees: child parameters
are computable without returned proof values, same-stratum ranks
strictly decrease, and lower-stratum calls follow a displayed order.
The tree is constructed with a primitive recursive size bound and
evaluated bottom-up.  This replaces earlier proofs that used induction
over an unbounded family of environments.  The fixed logical and
arithmetic leaves are cited as finite derivation recipes; validity is
not inferred from the truth or consistency of PA.

Third, the argument distinguishes current Foundation proof coding,
historical ordinary PAeq, and the numerical calculus.  It treats the
changed finite Foundation axiom and the metatheory-base comparison,
not merely the object-theory axiom translation.  Endpoint identity at
the corpus snapshot was checked explicitly: the two relevant endpoint
files match the baseline manifests and the Foundation gitlink remains
\path{32e1a0956a8622fad067328ca1959729a7634428}.  The structured
receipt is \path{proof/endpoint_identity.json}.

Finite checks were also rerun.  The source checks cover every literal
instruction and edge, the complete symbolic AX effect, and the exact
arithmetic templates.  The closure-constructor replay checks 308 closure
chains, 140 canonical induction readbacks, 36 existential projections,
four changed-leaf adapters and 8,828 distinct primitive proof nodes.
These checks support the concrete bindings and the displayed uniform
proofs; they are not a kernel verification of the complete arithmetic
derivation.  In particular, this work has not expanded every enormous
fixed numerical proof or implemented the complete PA compiler anew.

The retained archive is therefore an ordinary mathematical proof
dependency, with its exact sources and finite derivation recipes
available for further review.  The new relative proof changes only
the command implementation and the explicit generalization word, then
transfers through the same calibrated calculus.  Its conclusion should
be reported as a written iPRA equivalence with finite implementation
certificates.  The separate Lean formalization remains incomplete.
The longer audit and dependency map are in \path{proof/audit.md}.
